\documentclass[11pt,a4paper]{article}

\usepackage[latin1]{inputenc}
\usepackage[cyr]{aeguill}
\usepackage[francais]{babel}

% Les packages
%\usepackage[french]{babel}
\usepackage{fancybox}
\usepackage{graphics}
\usepackage{color}
\usepackage{latexsym}
\usepackage{amsmath,amsthm,amssymb}
\usepackage[plain]{fullpage} 
\usepackage{verbatim}
\usepackage{times,epsfig}
\usepackage{listings}
\usepackage{multirow}
\usepackage{amsmath,amsthm,amssymb,amscd,txfonts,bbm} % RAJOUT ICI !!!
\usepackage{graphics,epsfig} % RAJOUT ICI !!!
%\usepackage{algorithmic}
\usepackage[final]{pdfpages} 
\usepackage{fancyvrb}
\usepackage{enumerate}
\usepackage{listingsutf8}
\usepackage{comment}
\usepackage{minted}

% small et verbatim: to be coherent with \file
\newenvironment{smallverbatim}%
{\endgraf\small\verbatim}{\endverbatim}

\newenvironment{qv}
{\quote\Verbatim}
{\endquote\endVerbatim}

\usepackage{lastpage}

\newcommand{\terme}[1]{\texttt{#1}}
\newcommand{\termSet}[1]{\{\texttt{#1}\}}
\newcommand{\guill}[1]{<<\,#1\,>>}
% Macro pour l'inclusion de figure  
\newcommand{\fig}[2]{%
  \begin{figure*}[hbt]
   \begin{center}
      \includegraphics[width=6cm]{#1}
  \caption{\label{#1}  #2}
  \end{center}
 \end{figure*}
}
\newcommand{\figTex}[2]{%
  \begin{figure*}[hbt]
 \centerline{\input{#1.pstex_t}}
  \caption{\label{#1}  #2}
 \end{figure*}
}

% Macro commentaires
\definecolor{turquoise}{rgb}{0.20 ,0.60, 0.60}
\definecolor{vert}{rgb}{0.40,  0.60 ,0.00}
\definecolor{violet}{rgb}{0.80 , 0.00 , 0.80}
\definecolor{bleu}{rgb}{0.20,  0.20,  0.80}
\definecolor{rose}{rgb}{1.00,  0.20,  0.40}

\lstset{basicstyle=\footnotesize,showstringspaces=false,language=XML,breaklines=false,columns=fullflexible,frame=shadowbox, captionpos=b}

\newenvironment{solution}[0]{
~\\ \color{red} \textbf{Solution : }\color{black}}{%
}
\excludecomment{solution}


\parindent=0pt           
                             
\usepackage{url}

\newtheorem{Exercice}{Definition}


\begin{document}

\title{Sujet d'examen UE NFE204 - FOD Session 1}
\author{Année universitaire 2014-2015}
\date{Examen première session: 15 juin 2015}

\maketitle

\begin{center}
{\bf \Large Consignes}\\
\begin{tabular}{| p{15cm} |}
\hline
Dur\'ee de l'examen : 3h00; Documents non autoris\'es;  Calculatrice autoris\'ee\\

Les exercices sont indépendants les uns des autres. Merci de soigner votre écriture.\\
Ce sujet contient \pageref{LastPage} pages.\\

\textit{Important}\,: nous n'évaluons pas les réponses par la syntaxe. Ne vous
inquiétez pas d'utiliser un point-virgule à la place d'un point d'exclamation ou
autres détails mineurs
\\
\hline
\end{tabular}
\end{center}


\section*{Première partie : documents structurés (6 pts)}

Une application s'abonne à un flux de nouvelles dont voici un échantillon. 

\begin{minted}[frame=leftline,
               baselinestretch=.9,
               fontsize=\footnotesize,
               framesep=2mm]{xml}
<?xml version="1.0" encoding="UTF-8" ?>

<myrss version="2.0">
	<channel>
		<item>
			<title>Angela et nous</title>
			<description>Un nouveau sommet entre la France et l'Allemagne
				est prévu en Allemagne la semaine prochaine pour relancer l'UE.
			</description>
			<links>
				<link>lemonde.fr</link>
				<link>lesechos.fr</link>
			</links>
		</item>
		<item>
			<title>L'Union en crise</title>
			<description> L'Allemagne, la France, l'Italie doivent à nouveau se
				réunir pour étudier les demandes de la Grèce.
			</description>
			<links>
				<link>lefigaro.fr</link>
			</links>
		</item>
		<item>
			<title>Alexis à l'action</title>
			<description>Le nouveau premier ministre de la Grèce
				a prononcé son discours d'investiture.
			</description>
			<links>
				<link>libe.fr</link>
			</links>
		</item>
		<item>
			<title>Nos cousins transalpins</title>
			<description>La France et l'Italie partagent plus qu'une frontière
				commune: le point de vue de l'Italie.
			</description>
			<links>
				<link>courrier.org</link>
			</links>
		</item>
	</channel>
</myrss>
\end{minted}

Chaque nouvelle (\texttt{item}) résume donc un sujet et propose des liens
vers des médias où le sujet est développé. Les quatre éléments \texttt{item} seront désignés respectivement par
d1, d2, d3 et d4.

\begin{enumerate}
  \item Donnez la forme arborescente de ce document (ne recopiez
     pas tous les textes: la structure du document suffit).
  \item Proposez un format JSON pour représenter le même contenu (idem: la structure suffit).
  \item Dans une base relationnelle, comment modéliser
     l'information contenue dans ce document?
   \item Quelle est l'expression XPath pour obtenir tous les éléments \texttt{link}?
   \item Quelle est l'expression XPath pour obtenir
      les titres des items dont l'un des \texttt{link} est \texttt{lemonde.fr}?
\end{enumerate}

\section*{Deuxième partie : recherche d'information (6 pts)}

On veut indexer les nouvelles reçues de manière à pouvoir les rechercher
en fonction d'un pays. Le vocabulaire auquel on se restreint est donc celui
des noms de pays (Allemagne, France, Grèce, Italie).

\smallskip

\begin{enumerate}
   \item Donnez une matrice d'incidence contenant les tf pour les quatre pays,
      et un tableau donnant les idf (sans appliquer le $\log$). Mettez les noms de pays
      en ligne, et les documents en colonne.

\begin{solution}
\begin{center}
\begin{tabular}{l|c|c|c|c}
            & \textbf{d1} & \textbf{d2} & \textbf{d3} &\textbf{d4} \\ \hline  
   Allemagne  (2)   & 2          & 1           & 0           & 0        \\
   France    (4/3)  & 1           & 1           & 0          & 1         \\
   Grèce    (2)  & 0           & 1          & 1           & 0         \\
   Italie    (2)  & 0           & 1           & 0          & 2         \\ \hline
   \end{tabular}
\end{center}
\end{solution}

  \item Donner les résultats classés par similarité cosinus basée sur les tf 
  (on ignore l'idf) pour les requêtes suivantes. Expliquez brièvement le classement.
  
  \begin{itemize}
    \item Italie;
    \item Allemagne et France ;
    \item France et Grèce.
  \end{itemize}
\begin{solution}
Les normes
\begin{itemize}
  \item $||d_1|| = \sqrt{4+ 1} = \sqrt{5}$; $||d_2|| = \sqrt{1+1+1+1} = 2$; 
     $||d_3|| = \sqrt{1}$; 
    $||d_4|| = \sqrt{1+4}=\sqrt{5}$.
\end{itemize}

Les cosinus (requête non normalisée).

\begin{itemize}
  \item Italie: d2 : $\frac{1}{2}=0,5$;  
       d4 : $\frac{2}{\sqrt{5}} \simeq 0,89$ 
       
       d4 est premier car il mentionne deux fois l'Italie.
  
  \item Allemagne et France 
  
    d1 : $\frac{2+1}{\sqrt{5}}$; 
    d2 : $\frac{1+1}{2}$;
     d4 : $\frac{1}{\sqrt{5}}$; 
    
    d1 est premier comme on pouvait s'y attendre: il parle exclusivement
    de la France et de l'Allemagne. Viennent ensuite d2 puis d4.
    
  \item France et Grèce.
  
    d1 : $\frac{1}{\sqrt{5}}$; 
    d2 : $\frac{2}{2}$;
    d3 :  $\frac{1}{1}$  
    d4 :  $\frac{1}{\sqrt{5}}$  
    
  d2 et d3 arrivent à égalité. Intuitivement, il parlent tous les deux "à moitié"
   de France et de Grèce.  d1 et d4 parlent de France \emph{ou} de Grèce et aussi des autres
    pays.
 \end{itemize}

\end{solution}  
    
    \item Reprenons la dernière requête, "France et Grèce" et les documents d2 et d3.
       Qu'est-ce que la prise en compte de l'idf changerait au classement?
    
\begin{solution}
La Grèce a un idf plus élevé: d3 serait donc considéré come plus pertinent, 
car plus spécifique au besoin exprimé par la requête.
\end{solution}

  \end{enumerate}

\section*{Troisième partie : MapReduce (4 pts)}

Voici quelques analyses à exprimer avec MapReduce et Pig.

\begin{enumerate}
  \item   On veut  compter le nombre de nouvelles consacrées
à la Grèce publiées par chaque média. Décrivez le programme MapReduce
(fonctions de Map et de Reduce) qui produit le résultat souhaité. 
Donnez le pseudo-code de chaque
fonction, ou indiquez par un texte clair son déroulement.
Vous disposez d'une
fonction \textit{contains(\$texte, \$mot)} qui renvoie vrai
si \texttt{\$texte} contient \texttt{\$mot}. 

  \item Quel est le programme Pig qui exprime ce traitement MapReduce?
  
\end{enumerate}


\begin{solution}
\begin{enumerate}
 \item La fonction de Map
 
\begin{minted}{bash}
function map(doc)

  if (contains(doc.description, "Grèce") then
     for (link in doc.links) do
       emit (link, 1)
     done
  fi  
\end{minted}

La fonction de reduce, standard.

\begin{minted}{bash}
function reduce(link, counts)
  return [link, sum(counts)]
\end{minted}

\item

\begin{minted}{javascript}
greeknews = filter news by contains (description, "Grèce");
greeklinks = foreach greeknews generate description,
               flatten(greeknews.links) as link;
groupedlinks = group greeklinks by link;
result = foreach groupedlinks generate link, count(description)
\end{minted}
\end{enumerate}
\end{solution}

\section*{Quatrième partie : questions de cours (4 pts)}


\begin{enumerate}
  \item Expliquez les rôles de la réplication dans un système distribué. Quel est le
     plus important?
     
  \item Dans un traitement MapReduce, à quoi correspond l'opération dite de \textit{shuffle}?
  
  \item Expliquez ce qui caractérise un système "analytique" d'une part, un système
  "temps réel" d'autre part. Donnez des exemples de chaque type de système. 
      
   \item Qu'est-ce qui peut amener à choisir XML ou JSON pour la représentation de documents
      structurés? Citez, pour chaque format, 2 ou 3 spécificités.
\end{enumerate}

\end{document}
